% chapters/02/section-01.tex
\section{Permainan sebagai Sistem Waktu Nyata}
\label{sec:permainan-sistem}

Sebelum menulis satu baris kode pun, ada gunanya menyepakati apa yang
sebenarnya sedang dibangun. Bab ini dimulai dari pertanyaan paling dasar:
apa yang membuat sesuatu disebut permainan, dan apa yang membedakan program
permainan dari program lain yang juga berinteraksi dengan penggunanya.

\subsection{Unsur Pembentuk Permainan}

Dalam kajian desain permainan, sebuah permainan dipandang sebagai sistem
buatan tempat pemain terlibat dalam konflik yang diatur oleh aturan dan
menghasilkan keluaran yang dapat diukur. Rumusan itu memuat empat unsur yang
selalu hadir, bahkan pada permainan yang paling sederhana.

\begin{description}
  \item[Pemain] Pihak yang mengambil keputusan. Tanpa pemain, yang ada hanyalah
        simulasi yang berjalan sendiri.
  \item[Aturan] Batas tindakan yang boleh dilakukan pemain dan akibat dari
        setiap tindakan. Aturan inilah yang membedakan permainan dari
        sekadar mainan bebas.
  \item[Konflik atau tantangan] Hambatan yang membuat tujuan tidak dapat
        dicapai secara langsung, misalnya rintangan, lawan, atau keterbatasan
        sumber daya.
  \item[Keluaran terukur] Hasil yang dapat dinyatakan, misalnya menang atau
        kalah, skor, atau waktu penyelesaian.
\end{description}

Keempat unsur tersebut menjadi dasar pertanyaan desain yang akan Anda hadapi
berulang kali: tindakan apa yang boleh diambil pemain, apa akibatnya, dan
bagaimana keberhasilan dinyatakan. Pada tataran yang lebih teknis, keputusan
desain semacam itu diterjemahkan menjadi mekanik permainan, yaitu aturan yang
dijalankan oleh program \cite{unity-leveldesign}.

\subsection{Permainan Digital sebagai Sistem Waktu Nyata}

Yang membedakan program permainan dari program lain bukanlah grafisnya,
melainkan tuntutan waktu. Program pengolah data berjalan menurut urutan
perintah: membaca, menghitung, menulis, selesai. Program permainan tidak
memiliki akhir yang wajar seperti itu. Ia berjalan terus, bereaksi terhadap
masukan pemain dengan segera, dan menyajikan keadaan dunia yang terus berubah.

Tuntutan ``segera'' inilah yang dalam ilmu komputer disebut sistem waktu
nyata (\istilah{real-time system}): batas waktu tanggapan ditentukan oleh
kenyamanan orang yang memakainya, bukan oleh kecepatan mesin. Apabila
permainan membalas masukan setelah setengah detik, pemain merasakannya sebagai
kendali yang lambat dan tidak nyaman. Karena itu program permainan harus
menyelesaikan seluruh pekerjaannya dalam jangka waktu yang cukup pendek dan
mengulanginya terus-menerus.

Struktur pengulangan itulah yang disebut siklus permainan
(\istilah{game loop}), yaitu rangkaian tiga tahap---membaca masukan,
memperbarui keadaan, dan menayangkan hasilnya---yang dijalankan berulang
selama permainan berlangsung \cite{nystrom-gameloop}. Siklus tersebut menjadi
kerangka kerja seluruh bab selanjutnya, dan pada Bab~IV Anda akan menulisnya
sendiri baris demi baris.

\subsection{Subsistem Pembentuk Program Permainan}

Meskipun permainan terlihat sebagai satu kesatuan, programnya tersusun dari
beberapa subsistem yang tugasnya berbeda. Memisahkan pekerjaan menurut
subsistem sejak awal membuat program lebih mudah diperiksa dan diperbaiki
ketika jumlah kodenya bertambah. Tabel~\ref{tab:subsistem} memperlihatkan
pembagian tersebut beserta bab yang membahasnya.

\begin{table}[htbp]
\centering
\caption{Subsistem program permainan dan bab yang membahasnya}
\label{tab:subsistem}
{\setlength{\tabcolsep}{4pt}\footnotesize%
\begin{tabularx}{\textwidth}{@{}L{3.0cm} L{4.4cm} Y@{}}
\toprule
\headrow \textbf{Subsistem} & \textbf{Pekerjaan} & \textbf{Dibahas pada} \\
\midrule
Masukan &
Membaca peristiwa dari papan tombol, tetikus, atau peranti lain, lalu
menerjemahkannya menjadi tindakan yang bermakna. &
Bab~IV (peristiwa), Bab~V (papan tombol), Bab~XIII (kontrol pemain di Unity). \\
\addlinespace
Pengendali waktu &
Menjaga agar laju permainan konsisten, mengukur selang waktu antar-bingkai, dan
menentukan panjang langkah simulasi. &
Bab~IV (pembatasan laju bingkai), Bab~V (\istilah{delta time}). \\
\addlinespace
Simulasi &
Memperbarui keadaan dunia: posisi, kecepatan, keadaan NPC, skor, dan fase
permainan. &
Bab~V--VI (gerak dan fisika), Bab~VII (tabrakan), Bab~VIII--XI (pengendalian
NPC). \\
\addlinespace
Aturan dan tabrakan &
Menentukan kapan dua objek bersentuhan dan apa akibatnya bagi keadaan
permainan. &
Bab~VII (deteksi tabrakan), Bab~XIII (fisika mesin). \\
\addlinespace
Penyajian &
Mengubah keadaan permainan menjadi gambar di layar. &
Bab~IV (jendela dan penayangan), Bab~XIV (aset visual, animasi, antarmuka). \\
\addlinespace
Audio &
Membangkitkan bunyi sebagai tanggapan atas peristiwa permainan. &
Bab~XV. \\
\addlinespace
Data dan penyimpanan &
Menyimpan peta, pengaturan, dan kemajuan permainan. &
Bab~IX--X (data peta), Bab~XIV (aset), Bab~XVI (berkas terpasang). \\
\bottomrule
\end{tabularx}}
\tabsource{Sumber: disusun dari RPS mata kuliah \NamaMK\ dan sumber \mbox{[R1]},
\mbox{[R5]}, \mbox{[R12]}.}
\end{table}

\subsection{Peta Konsep Bab Ini}

Gambar~\ref{fig:peta-konsep} memetakan keenam tema landasan teori pada bab ini
beserta bab yang menerapkannya. Gambar tersebut dapat dipakai sebagai
penunjuk arah: apabila pada bab praktik Anda menemui istilah yang terasa asing,
tempat pertama yang sebaiknya dibuka adalah bagian teori yang bersangkutan di
bab ini.

\begin{figure}[htbp]
\centering
% chapters/02/figures/peta-konsep.tikz
% Peta konsep Bab II: setiap tema landasan teori dan bab yang menerapkannya.
\begin{tikzpicture}[
  konsep/.style={
    rectangle, rounded corners=2pt, draw=ubPrimary, thick, fill=ubAlt,
    align=left, text width=56mm, inner sep=4pt, minimum height=9mm,
    font=\footnotesize
  },
  bab/.style={
    rectangle, rounded corners=2pt, draw=ubAccent, thick, fill=white,
    align=left, text width=42mm, inner sep=4pt, minimum height=9mm,
    font=\footnotesize
  },
  panah/.style={-{Stealth[length=2.2mm]}, thick, draw=ubAccent}
]

\node[konsep] (k1) at (0,0)
  {\textbf{2.1} Permainan sebagai sistem waktu nyata};
\node[konsep] (k2) at (0,-1.35)
  {\textbf{2.2} Model waktu: bingkai dan langkah simulasi};
\node[konsep] (k3) at (0,-2.70)
  {\textbf{2.3} Ruang, besaran, dan representasi keadaan};
\node[konsep] (k4) at (0,-4.05)
  {\textbf{2.4} Simulasi fisika diskret dan integrasi numerik};
\node[konsep] (k5) at (0,-5.40)
  {\textbf{2.5} Keputusan dan navigasi pada NPC};
\node[konsep] (k6) at (0,-6.75)
  {\textbf{2.6} Pustaka dan mesin: dua cara membangun};

\node[bab] (b1) at (7.6,0)        {Bab III--IV};
\node[bab] (b2) at (7.6,-1.35)    {Bab IV--VI};
\node[bab] (b3) at (7.6,-2.70)    {Bab V, VII, XIII};
\node[bab] (b4) at (7.6,-4.05)    {Bab VI, XIII};
\node[bab] (b5) at (7.6,-5.40)    {Bab VIII--XI};
\node[bab] (b6) at (7.6,-6.75)    {Bab XII--XVII};

\foreach \i in {1,...,6} {
  \draw[panah] (k\i) -- (b\i);
}
\end{tikzpicture}

\caption{Peta konsep landasan teori dan bab yang menerapkannya}
\label{fig:peta-konsep}
\end{figure}

Tema \textbf{2.2} sampai \textbf{2.4} menopang Blok~I dan
Blok~III sekaligus, sedangkan tema \textbf{2.5} khusus menopang pengendalian
NPC pada Blok~II. Susunan itu sejalan dengan urutan mata kuliah: pemahaman
tentang waktu, ruang, dan fisika dipakai kembali ketika Anda bekerja dengan
mesin permainan, dengan tingkat tanggung jawab yang berbeda.
